% ============================================================
% Plantilla: Ensayo con normas MLA 9.ª edición
%   \autocite[45]{clave} -> (Autor 45)
% Recompilar ejecuta biber automáticamente.
% ============================================================
\documentclass[12pt,a4paper]{article}

\usepackage[utf8]{inputenc}
\usepackage[T1]{fontenc}
\usepackage{mathptmx}
\usepackage[spanish]{babel}
\usepackage[margin=2.54cm,headheight=15pt]{geometry}
\usepackage{setspace}
\usepackage{csquotes}
\usepackage[style=mla,backend=biber]{biblatex}
\usepackage[hidelinks]{hyperref}

\addbibresource{referencias.bib}
\doublespacing
\setlength{\parindent}{1.27cm}

% MLA: apellido del autor y número de página arriba a la derecha
\newcommand{\apellido}{Apellido}
\usepackage{fancyhdr}
\pagestyle{fancy}
\fancyhf{}
\fancyhead[R]{\apellido\ \thepage}
\renewcommand{\headrulewidth}{0pt}

\begin{document}

% ---------- Encabezado del estudiante (sin portada) ----------
\noindent Nombre Apellido\\
Nombre del Docente\\
Nombre de la Asignatura\\
\today

\begin{center}
  Título del Ensayo
\end{center}

El ensayo comienza aquí. En MLA las citas indican el autor y la página, por ejemplo
\autocite[12]{lamport1994}. Si el autor ya se menciona en el texto, como Knuth, basta
con la página \autocite[45]{knuth1984}.

Cada párrafo desarrolla una idea que apoya la tesis del ensayo y la respalda con
fuentes \autocite[3]{mla2021}.

En la conclusión se retoma la tesis y se resumen los argumentos principales
\autocite[102]{goossens2004}.

\newpage
\printbibliography[title=Obras citadas]

\end{document}
